\section{Model}
\label{sec:model}

The link is represented by three elements in series: the rectifier station, the
DC line, and the inverter station. Each is a single algebraic relation.

\subsection{Line resistance}
The resistance of one pole conductor over the whole route is the specific
resistance of a single subconductor, divided by the number of subconductors in
the bundle (they are electrically in parallel), converted to $\Omega$/mi and
multiplied by the length:
\begin{equation}
  R_{HV} = \frac{r'_{HV}\cdot 5.28}{n_{HV}}\; L ,
  \qquad
  R_{DMR} = \frac{r'_{DMR}\cdot 5.28}{n_{DMR}}\; L .
  \label{eq:resistance}
\end{equation}
The resistance that matters for the flow is that of the complete circuit ---
out and back --- which depends on the configuration:
\begin{equation}
  R_{loop} =
  \begin{cases}
    2\,R_{HV}          & \text{symmetrical monopole, and balanced bipole,}\\[2pt]
    R_{HV} + R_{DMR}   & \text{bipole with one pole out (DMR return).}
  \end{cases}
  \label{eq:loop}
\end{equation}

Equation~\eqref{eq:resistance} is a DC resistance at the temperature of the
underlying conductor table. No temperature correction, skin effect or
proximity effect is applied; on a DC line the first is the only one that
matters and it is handled by choosing the conductor data at the design
temperature (Section~\ref{sec:assumptions}).

\subsection{Converter stations}
Station losses are taken as a fraction of the power passing through the
station. The two stations therefore differ only in which side of the
conversion is known:
\begin{align}
  \text{rectifier (AC}\rightarrow\text{DC):} \quad
    & P_{dc,rec} = P_{ac,rec}\,\bigl(1-p_{rec}\bigr)
      \quad\Longleftrightarrow\quad
      P_{ac,rec} = \frac{P_{dc,rec}}{1-p_{rec}},
      \label{eq:rect}\\[2pt]
  \text{inverter (DC}\rightarrow\text{AC):} \quad
    & P_{ac,inv} = P_{dc,inv}\,\bigl(1-p_{inv}\bigr)
      \quad\Longleftrightarrow\quad
      P_{dc,inv} = \frac{P_{ac,inv}}{1-p_{inv}}.
      \label{eq:inv}
\end{align}
A single fraction lumps everything inside the station fence: valve conduction
and switching losses, converter transformer and reactor losses, cooling and
auxiliary supply. The loss fraction is a straight percentage of throughput, so
it scales linearly with power and does not capture the no-load component of a
real loss curve; see Section~\ref{sec:assumptions}.

Because power is specified at the receiving AC terminals,
Equation~\eqref{eq:inv} is the \emph{entry point} of the whole calculation: it
converts the given $P_{ac,inv}$ into the DC power that must arrive at the
inverter DC terminals.

\subsection{DC line}
With the DC voltage held at $V_{d}$ pole-to-ground at the rectifier, the power
leaving the rectifier and the resistive loss of the line are
\begin{equation}
  P_{dc,rec} = k\,V_{d}\,I_{d},
  \qquad
  P_{loss,line} = R_{loop}\,I_{d}^{\,2},
  \qquad
  k =
  \begin{cases}
    2 & \text{two poles energised,}\\
    1 & \text{one pole out,}
  \end{cases}
  \label{eq:dcline}
\end{equation}
and the power arriving at the inverter DC terminals is the difference:
\begin{equation}
  P_{dc,inv} = k\,V_{d}\,I_{d} - R_{loop}\,I_{d}^{\,2}.
  \label{eq:flow}
\end{equation}
Equation~\eqref{eq:flow} is the load flow equation of the link. Everything in
Section~\ref{sec:solution} follows from solving it.

The loss splits between the conductors in proportion to their resistances,
which is reported separately because the DMR contribution is what distinguishes
the outage case:
\begin{equation}
  P_{loss,HV} =
  \begin{cases}
    2R_{HV} I_{d}^{\,2} & \text{balanced,}\\
    R_{HV} I_{d}^{\,2}  & \text{one pole out,}
  \end{cases}
  \qquad
  P_{loss,DMR} =
  \begin{cases}
    0                    & \text{balanced,}\\
    R_{DMR} I_{d}^{\,2}  & \text{one pole out.}
  \end{cases}
  \label{eq:losssplit}
\end{equation}

\subsection{Receiving-end voltage}
The voltage at the far end of the line is the sending-end voltage less the drop
along the conductors the current actually traverses. With two poles energised
the arrangement is symmetric --- each conductor drops $I_{d}R_{HV}$ and the
midpoint stays at ground --- so the pole-to-ground voltage at the inverter is
\begin{equation}
  V_{d,inv} = V_{d} - I_{d}\,R_{HV}
  \qquad\text{(two poles energised),}
  \label{eq:vinv}
\end{equation}
which is consistent with Equation~\eqref{eq:flow}, since
$2V_{d,inv}I_{d} = 2V_{d}I_{d} - 2R_{HV}I_{d}^{2}$. With one pole out the whole
loop drop appears between the healthy pole and the DMR, so the corresponding
figure is a pole-to-DMR voltage:
\begin{equation}
  V_{d,inv} = V_{d} - I_{d}\,R_{loop}
  \qquad\text{(one pole out, pole-to-DMR).}
  \label{eq:vinvmono}
\end{equation}
